% Self-contained section snippet. Requires amsmath,amssymb,amsthm.
% Theorem environment names used: theorem,lemma,proof,remark.
\section{A multivariate exact selector without iterated ramification}
\label{sec:multi}
The preceding fixed-field results lead to a new question: must a successive
choice of torsion coordinates enlarge the coefficient field at every step?
The recurrence in the earlier report does so and can grow as an iterated
exponential. A finite list of multiplicative curves gives a different
reduction. The curve is selected before choosing its torsion parameter;
all its defining coefficients are integers. The coefficient field therefore
stays fixed throughout the reduction.

Write $v(p)=1$ and, for an integral polynomial $F$,
\[
 G_\rho(F)=\min_{\alpha}\{v(c_\alpha)+\rho|\alpha|\}.
\]
The same notation with several subscripts denotes different positive
weights on the coordinates. These Gauss valuations are multiplicative
and extend to valuations of the rational-function field.

\begin{lemma}[Finite slopes preserve individual bounds]
\label{lem:finite-slopes}
Let $F_1,\ldots,F_t$ be nonzero integral polynomials in
$X,Y,Z_1,\ldots,Z_q$, over any valued field extending $\mathbf Q_p$.
Fix $\rho>0$, put $A_i=G_\rho(F_i)$, and assume $\sum_i A_i\le B$.
Set
\[
 M=\left\lfloor\frac B\rho\right\rfloor+1,
 \qquad r=\frac\rho{M^2},\qquad b_n(T)=(1+T)^n-1.
\]
For some $1\le n\le M$, simultaneously for every $i$,
\[
 G_{r,\rho,\ldots,\rho}
   \bigl(F_i(T,b_n(T),Z_1,\ldots,Z_q)\bigr)\le A_i.
\]
In particular, the same holds after lowering all remaining weights to $r$.
No bound on the degrees or on $t$ is required.
\end{lemma}

\begin{proof}
Consider one $F$, with $A=G_\rho(F)$, and put
$k=\lfloor A/\rho\rfloor+1$. We prove that fewer than $k$ slopes are bad,
where bad means that the restricted Gauss valuation is strictly greater
than $A$, including $+\infty$.
Suppose a set $J\subseteq\{1,\ldots,M\}$ of $k$ slopes were bad. Divide in
$Y$ by the monic polynomial
\[
 Q_J(X,Y)=\prod_{n\in J}(Y-b_n(X)).
\]
The quotient and remainder are integral; write $F=Q_JS+R$ with
$\deg_Y R<k$. Each factor of $Q_J$ has Gauss valuation exactly $\rho$,
so $G_\rho(Q_JS)\ge k\rho>A$ and $G_\rho(R)=A$.

Choose a monomial attaining the last minimum. If its degree in $X,Y$ is
zero, then $G_\rho(R(0,0,Z))\le A$. But
$R(0,0,Z)=F(0,0,Z)$, the constant coefficient in $T$ of every restriction.
This makes every slope good, a contradiction. Otherwise its degree in
$X,Y$ is at least one, whence
\begin{equation}\label{eq:slope-saving}
 G_{r,r,\rho,\ldots,\rho}(R)\le A-(\rho-r).
\end{equation}

Lagrange interpolation gives
\[
 R(X,Y,Z)=\sum_{n\in J}F(X,b_n(X),Z)
   \frac{\prod_{m\in J\setminus\{n\}}(Y-b_m(X))}{D_n(X)},
 \quad
 D_n(X)=\prod_{m\in J\setminus\{n\}}(b_n(X)-b_m(X)).
\]
The numerator product has Gauss valuation $(k-1)r$. For a positive
integer $d$, the coefficient of $T^{p^{v_p(d)}}$ in
$(1+T)^d-1$ is a $p$-adic unit: reduce
$(1+T)^d=(1+T^{p^{v_p(d)}})^{d/p^{v_p(d)}}$ modulo $p$.
Consequently
\[
 G_r(b_n-b_m)\le p^{v_p(|n-m|)}r.
\]
Define
\[
 C_J=\max_{n\in J}
 \sum_{m\in J\setminus\{n\}}\bigl(p^{v_p(|n-m|)}-1\bigr).
\]
Since all the slopes in $J$ were bad, the Gauss triangle inequality in
this interpolation identity yields
\begin{equation}\label{eq:slope-interpolation}
 G_{r,r,\rho,\ldots,\rho}(R)>A-rC_J.
\end{equation}
For $M\ge2$,
$C_J\le(k-1)(M-2)\le(M-1)(M-2)$; for $M=1$ it is zero.
Thus $C_J+1\le M^2$ and $rC_J\le\rho-r$.
Equations~\eqref{eq:slope-saving}--\eqref{eq:slope-interpolation}
contradict each other.

Each $F_i$ therefore has at most $\lfloor A_i/\rho\rfloor$ bad slopes.
Their union has size at most
$\sum_i\lfloor A_i/\rho\rfloor\le\lfloor B/\rho\rfloor=M-1$.
At least one of the $M$ slopes works for all $i$.
\end{proof}

\begin{theorem}[A bound with no ramification tower]
\label{thm:no-tower}
Let $K/\mathbf Q_p$ be finite of absolute ramification index $e$, put $w=v_p(e)$, and let
$F_i\in\mathcal O_K[T_1,\ldots,T_d]\setminus\{0\}$, $d\ge1$, be a finite
packet. Fix $\rho>0$, put $A_i=G_\rho(F_i)$, and assume
$\sum_i A_i\le B$. Let $h_0\ge1$ satisfy
$1/\varphi(p^{h_0})\le\rho$.
There are roots of unity
$\zeta_1,\ldots,\zeta_d$ of respective orders $p^{h_1},\ldots,p^{h_d}$,
with
\[
 h_0\le h_j\le H,\qquad
 v\bigl(F_i(\zeta_1-1,\ldots,\zeta_d-1)\bigr)\le A_i
 \quad\text{for every }i,
\]
where it suffices to take
\begin{equation}\label{eq:no-tower-bound}
 H=h_0+
 \left\lceil(3^{d-1}-1)\log_p\!\left(1+\frac B\rho\right)\right\rceil
 +\lfloor p^w B\rfloor.
\end{equation}
\end{theorem}

\begin{proof}
Put $\rho_0=\rho$ and, for $j=0,\ldots,d-2$, recursively define
\[
 M_j=\left\lfloor\frac B{\rho_j}\right\rfloor+1,
 \qquad \rho_{j+1}=\frac{\rho_j}{M_j^2}.
\]
At each stage apply Lemma~\ref{lem:finite-slopes} to the first remaining
parameter and the next variable, then lower the weights of the other
variables to $\rho_{j+1}$. Every individual Gauss valuation decreases or
stays the same. All coefficients remain in $\mathcal O_K$.
We obtain integers $a_1=1,a_2,\ldots,a_d>0$ and
\[
 P_i(T)=F_i((1+T)^{a_1}-1,\ldots,(1+T)^{a_d}-1),
 \qquad G_{\rho_{d-1}}(P_i)\le A_i.
\]
One can take $a_{j+2}\le M_j$. Write
$P=\prod_{j=0}^{d-2}M_j$, with empty product $1$. Then $a_j\le P$ and
$\rho_{d-1}=\rho/P^2$.

Put $\bar h=h_0+\lceil2\log_pP\rceil$. Its cyclotomic radius satisfies
$r_{\bar h}\le\rho/P^2$. The scalar fixed-field selector, begun at
$\bar h$, chooses a primitive $p^h$-th root $\zeta$ with
\[
 \bar h\le h\le\bar h+\lfloor p^w B\rfloor,
 \qquad v(P_i(\zeta-1))\le A_i \quad\text{for all }i.
\]
Apply Theorem~\ref{thm:scalar} with thresholds $A_i$: each blocked level costs at least $p^{-w}$, so the total cost of these $\lfloor p^w B\rfloor+1$ levels exceeds $B$. The thresholds dominate the final Gauss valuations; they need not be actual evaluations at a common witness.
Set $\zeta_j=\zeta^{a_j}$. Its order exponent is
$h-v_p(a_j)$, which is at least
\[
 h_0+\lceil2\log_pP\rceil-\lfloor\log_pP\rfloor\ge h_0.
\]
This proves the required coordinate condition.

For the explicit estimate, put $u_j=1+B/\rho_j$. Then
\[
 u_{j+1}=1+(u_j-1)M_j^2\le u_j^3,
 \qquad u_j\le u_0^{3^j}.
\]
It follows that
\[
 2\log_pP\le(3^{d-1}-1)\log_p(1+B/\rho).
\]
This gives~\eqref{eq:no-tower-bound}. The cases of one variable, $B=0$, and $M_j=1$ are included: the empty product is $1$ and the logarithmic increment is zero when appropriate.
\end{proof}

\paragraph{Bounded power series and arbitrary packets.}
The same bound applies to nonzero elements of
$\mathcal O_K[[T_1,\ldots,T_d]]$. Choose a truncation degree $D$ with
$(D+1)r_H>B$. The truncation preserves every initial Gauss minimum, and
its omitted tail has valuation $>B$ at every candidate torsion tuple.
The polynomial theorem therefore transfers with exact equality of the
selected evaluations' valuations. For an arbitrary family whose sum of
Gauss budgets, defined as the supremum of finite partial sums, is at most
$B$, apply this result to every finite subfamily. The candidate torsion
grid is finite and independent of the subfamily. Its nonempty good sets
have the finite intersection property, giving a common tuple for the
whole family.

\paragraph{Why the finite list is needed.}
One prescribed multiplicative curve cannot suffice: for coprime positive
integers $a,b$, the curve
$((1+T)^a-1,(1+T)^b-1)$ is annihilated by
\[
 (1+T_1)^b-(1+T_2)^a,
\]
an integral polynomial of Gauss valuation exactly $\rho$. At least one
linear coefficient is a $p$-adic unit. The finite-slope selection removes
this obstruction while the interpolation estimate, rather than mere
nonvanishing, preserves the individual valuation targets.

\paragraph{Scope.}
This settles the explicit local multivariate growth question posed in the
previous report. It preserves a fixed coefficient field before selecting
the order bound. It does not authorize replacing a theorem whose
coefficient ring is allowed to ramify after that bound is selected, and
does not construct arithmetic interpolation data or nonzero witnesses.
The dimension factor $3^{d-1}$ and slope range are not claimed optimal.
